\begin{table}[!ht]
\centering
\caption{H4 --- Do inflation expectations lead or follow? Net directional connectedness}
\label{tab:h6net}
\setlength{\tabcolsep}{6pt}\footnotesize
\begin{threeparttable}
\begin{tabular}{l c c c}
\toprule
Sample & $\tau=0.5$ & $\tau=0.7$ & $\tau=0.9$ \\
\midrule
Full (1960--2026) & +4.8 & +13.7 & +10.5 \\
Great Inflation (1967--82) & +0.5 & -2.2 & -20.2 \\
Core (1983--2019) & +3.0 & +8.8 & +29.1 \\
COVID (2020--26) & -5.5 & -4.1 & -6.7 \\
\bottomrule
\end{tabular}
\begin{tablenotes}[flushleft]\footnotesize
\item \textit{Notes:} Net directional connectedness of quarterly Michigan inflation expectations (\texttt{MICH\_QTR}, from 1960) versus overall CPI in a two-variable system, $n_{\text{lag}}=2$. $\text{NET}=\text{TO}-\text{FROM}$; $>0$ means expectations lead (transmit), $<0$ means they follow. Pooled, expectations lead in the upper tail --- but this reflects the anchored 1983--2019 core (NET $+29.1$ at $\tau=0.9$); conditional on a high-inflation episode the sign reverses (Great Inflation $-20.2$, COVID $-6.7$), consistent with adaptive expectations. The COVID window (25 quarters) is short.
\end{tablenotes}
\end{threeparttable}
\end{table}

